\section*{Exercises}
Exercises~\ref{ex:4.1}--\ref{ex:4.4} are core exercises, and Exercises~\ref{ex:4.5} and~\ref{ex:4.6} are transfer exercises. Complete solutions are in Appendix~\ref{app:sola}, under ``Week~4.''
\begin{exercise}\label{ex:4.1}Give a direct proof that if $a$ and $b$ are even integers, then $ab$ is divisible by four.\end{exercise}
\begin{exercise}\label{ex:4.2}Let $n$ be an integer. Prove, by proving the contrapositive, that if $3n+1$ is even, then $n$ is odd.\end{exercise}
\begin{exercise}\label{ex:4.3}Let $r$ be a rational number and $u$ an irrational number. Prove that $r+u$ is irrational. Your proof will need the fact that the difference of two rational numbers is rational; explain why this fact holds.\end{exercise}
\begin{exercise}\label{ex:4.4}Prove that the equation $x=3x-4$ has exactly one real solution. Give the existence part and the uniqueness part separately.\end{exercise}
\begin{exercise}\label{ex:4.5}An assistant is asked to solve the equation $x^2=2$ for a real number $x$, and answers ``so $x=\sqrt2$.'' What has the assistant missed? Find all real solutions of $x^2=2$, and prove that your list is complete.\end{exercise}
\begin{exercise}\label{ex:4.6}Prove that $\sqrt8$ is irrational in two ways. First, deduce it from Theorem~\ref{thm:sqrt2}. Second, without using Theorem~\ref{thm:sqrt2}, adapt its proof so that it works for $\sqrt8$, and explain how the parity steps have to change.\end{exercise}
